
\begin{frame}[fragile]{Notions élémentaires de logique formelle}

Cette session présente les notions essentielles de logique formelle à connaître: 

\vfill

\begin{redbullet}
  \item Le calcul propositionnel: valeurs de vérité et formules
  \item Les prédicats 
  \item Collections et quantificateurs
  \item Importance pour les bases de données et SQL
\end{redbullet}

\vfill
\begin{blocgauche}
\alert{Ces diapositives correspondent au support en ligne disponible sur
le site \url{http://sql.bdpedia.fr/}}
\end{blocgauche}

\end{frame}


%%%%%%ù

\begin{frame}[fragile]{Calcul propositionnel}

\alert{Proposition}: énoncé auquel on peut affecter deux valeurs de vérité, Vrai ou Faux.

\vfill

\alert{Formule}: expression basée sur des propositions et leur combinaison
par des \alert{connecteurs logiques}.


\vfill
\begin{redbullet}
   \item la conjonction, notée $\land$
   \item  la disjonction, notée $\lor$
   \item la négation, notée $\neg$
\end{redbullet}

\vfill

Exemple: $(p \land q) \lor r$
\end{frame}



\begin{frame}[fragile]{Valeur de vérité d'une formule}

Induite à partir de valeurs de vérité des propositions, et des règles suivantes:

\vfill

\begin{tabular}{c|c|c|c|c}
$p$ & $q$ & $p \land q$ & $p \lor q$ & $\neg p$\\ \hline
Vrai & Vrai & Vrai & Vrai &  Faux \\ \hline
Vrai & Faux & Faux & Vrai &  Faux \\ \hline
Faux & Vrai & Faux & Vrai &  Vrai \\ \hline
Faux & Faux & Faux & Faux &  Vrai \\ \hline
\end{tabular}
\vfill

Exemple: $(p \land q) \lor r$ est Vrai pour les valeurs V, V, F de $p, q, r$

\end{frame}



\begin{frame}[fragile]{Equivalences}

Un informaticien averti sait manier avec agilité les équivalences. Quelques exemples.

\vfill

\begin{redbullet}
   \item  $\neg (\neg F)$ est équivalente à $F$ 
  \item  $F \lor (F_1 \land F_2)$ est équivalente à $(F\lor F_1) \land (F \lor F_2)$ (distribution)
   \item  $F \land (F_1 \lor F_2)$ est équivalente à $(F\land F_1) \lor (F \land F_2)$ (distribution)
   \item  $\neg (F_1 \land F_2)$ est équivalente à $(\neg F_1) \lor (\neg F_2)$ (loi DeMorgan)
   \item  $\neg (F_1 \lor F_2)$ est équivalente à $(\neg F_1) \land (\neg F_2)$ (loi DeMorgan)
\end{redbullet}

\vfill

Donc $p \lor \neg (p \land \neg q)$ est une \alert{tautologie}.
\end{frame}


\begin{frame}{Prédicats}

Extension puissante des propositions : \alert{construire des énoncés sur des ``objets''}.

\vfill

Le prédicat $Compose (X, Y)$ permet de construire des énoncés de la forme:

\begin{redbullet}
   \item $Compose  (\text{'Mozart'}, \text{'Don Giovanni'})$
   \item $Compose  (\text{'Debussy'}, \text{'La mer'})$
  \item  Etc.
\end{redbullet}

\vfill

Ce sont des nuplets (ou des atômes, ou des faits). 

\begin{redbullet}
   \item Il en existe une infinité
   \item Un \alert{contexte} (ou ``interprétation'') définit ceux qui sont vrais / faux.
\end{redbullet}

\vfill

\begin{blocgauche}
\begin{remblock}{Base de données?}
C'est un contexte donnant un ensemble fini de faits vrais. Tous les autres sont considérés comme faux.
\end{remblock}
\end{blocgauche}

\end{frame}




\begin{frame}{Nuplets ouverts et fermés}

Un nuplet énoncé avec des constantes est un nuplet \alert{fermé}. 
$$Compose  (\text{'Mozart'}, \text{'Don Giovanni'})$$
\vfill
Un nuplet énoncé avec au moins une variable est un nuplet \alert{ouvert}. 
$$Compose  (X, \text{'Don Giovanni'})$$
\vfill
Un nuplet ouvert désigne une infinité de faits possibles. 

\vfill


\begin{blocgauche}
\begin{remblock}{Interêt?}
En général on s'intéresse aux valeurs de $X$ pour lesquelles les faits sont vrais.
\alert{On a effectué une requête}.
\end{remblock}
\end{blocgauche}


\end{frame}



\begin{frame}{Collections et quantificateurs}

Les nuplets libres expriment des contraintes sur un fait. 

\vfill


On peut exprimer des contraintes sur des collections de faits avec les quantificateurs.


\begin{redbullet}
   \item $\exists x P(x)$ est vraie s'il existe \alert{au moins} une affectation de $x$
     pour laquelle $P(x)$ est 
     vraie.
   \item $\forall x P(x)$ est vraie si $P(x)$ est 
     vraie pour \alert{toutes} les valeurs de $x$.
\end{redbullet}

\vfill

\begin{blocgauche}
\alert{Variables libres et liées}: un variable quantifiée est libre; sinon elle est liée.
\end{blocgauche}

\end{frame}


\begin{frame}[fragile]{SQL = formules logique}

\alert{Requête SQL } = une formule avec des variables libres.

\vfill

\alert{Résultat d'une requête} = les valeurs des variables libres qui satisfont la formule.

\vfill
 La formule $Compose  (X, \text{'Don Giovanni'})$ s'écrit en SQL

\begin{minted}[baselinestretch=.9,
               fontsize=\small,
               framesep=2mm]{sql}
     select compositeur
       from Compose
       where oeuvre='Don Giovanni'
\end{minted}

SQL c'est une syntaxe pour écrire des formules.
\vfill


\begin{blocgauche}
\begin{remblock}{Déclarativité}
On ne dit pas comment on calcule!
\end{remblock}
\end{blocgauche}

\end{frame}


\begin{frame}[fragile]{Exemples}

Deux relations/prédicats:  

\begin{redbullet}
  \item \texttt{Expert (id\_expert, nom)} 
  \item \texttt{Manuscrit (id\_manuscrit, auteur, titre, 
id\_expert, commentaire)}
\end{redbullet}

\vfill 

$Manuscrit (\_, \text{'Proust'}, t, x, \_) \land Expert (x, n)$

\begin{minted}[baselinestretch=.9,
               fontsize=\small,
               framesep=2mm]{sql}
select titre, nom from Expert, Manuscrit
where Expert.id_expert = Manuscrit.id_expert
and auteur = 'Proust'
\end{minted}

\vfill

$Expert (x, n) \land \exists Manuscrit (\_, \text{'Proust'}, \_, x, \_) $

\begin{minted}[baselinestretch=.9,
               fontsize=\small,
               framesep=2mm]{sql}
select nom from Expert as e 
where exists (select *  from Manuscrit as m
              where e.id_expert = m.id_expert
              and auteur = 'Proust')
\end{minted}

\end{frame}


\begin{frame}{À retenir}

SQL est un langage \alert{déclaratif} qui exprime par une formule logique
les propriétés du résultat à construire. 

\vfill
Avantages prouvés et éprouvés depuis les années 1970.

\begin{redbullet}
  \item  signification précise, non ambigue
  \item algorithmes efficaces 
  \item  langage robuste, universellement connu et adopté, 
     \alert{normalisé}
  \item \alert{déclarativité}: SQL 
    ne donne  aucune indication sur la manière dont le système doit trouver le résultat.
\end{redbullet}

\end{frame}

